% CKPT-003 theory claim ledger. This is intentionally not a completed 20-page report.
\section{CKPT-003 theory programme and claim ledger}
\label{sec:ckpt003-theory-scope}

\textbf{Status.} This document is a proof-plan and scope ledger for
\texttt{v1.0.0-scale}; it contains no completed generalization theorem,
optimization theorem, sparse-memory capacity theorem, or independent review.
A statement is not promoted from ``open'' merely because a related experiment
has run.

\begin{center}
\begin{tabular}{p{0.22\linewidth}p{0.29\linewidth}p{0.17\linewidth}p{0.22\linewidth}}
\toprule
Workstream & Required formal object & Status & Required boundary \\
\midrule
Generalization & VC/pseudodimension or Rademacher bound for a defined function class & Open & Specify activation, norm/rank constraints, data distribution and probability. \\
Gradient flow & Dynamics for the declared low-rank parameterization & Open & Separate continuous-time surrogate from discrete optimizer. \\
Saddles & Stationary-point classification under named loss/constraints & Open & Do not infer global convergence. \\
Memory information & Capacity/encoding theorem for a stated memory channel & Open & Separate ideal address space from implemented hash/cache. \\
Addressability & Map from keys to physical slots and retrieval success & Open & Include collisions, eviction and decay. \\
Effective dimension & Generalization relation under stated data/model assumptions & Open & Do not equate parameter count with effective dimension. \\
\bottomrule
\end{tabular}
\end{center}

\subsection{Required proof structure}
Each prospective result must state definitions, assumptions, theorem, proof,
edge cases, comparison with the implementation, and falsifiable experimental
consequence. Results for an idealized random hash or an activation-free
Kronecker chain may be useful reference models, but may not be described as a
proof about the full trained system without an explicit reduction.

\subsection{Experimental interface}
For each proposed theoretical quantity, record its estimator, confidence
interval, finite-sample bias, architecture/configuration, data split and code
revision. The 1M-context stability gate measures only loss drift; it is not an
estimator of VC dimension, Rademacher complexity, information capacity or
optimization convergence.

\subsection{Review gate}
A reviewer-ready report must be at least twenty substantive pages of complete
claims, derivations, references and limitations---not padded prose---and must
receive independent mathematical review. This outline does not satisfy that
gate.
